\def\stat{zats+kor}



\def\tit{ИНФОРМАЦИОННАЯ БЕЗОПАСНОСТЬ СИТУАЦИОННЫХ 

ЦЕНТРОВ$^*$}



\def\titkol{Информационная безопасность ситуационных 

центров}



\def\autkol{А.\,А.~Зацаринный, В.\,И.~Королёв}



\def\aut{А.\,А.~Зацаринный$^1$, В.\,И.~Королёв$^2$}



\titel{\tit}{\aut}{\autkol}{\titkol}



\index{Зацаринный А.\,А.}

\index{Королёв В.\,И.}

\index{Zatsarinny A.\,A.}

\index{Korolev V.\,I.}



{\renewcommand{\thefootnote}{\fnsymbol{footnote}}

\footnotetext[1]

{Работа выполнена при поддержке РФФИ (проект 15-29-07981 офи-м).}}



\renewcommand{\thefootnote}{\arabic{footnote}}

\footnotetext[1]{Федеральный исследовательский

  центр <<Информатика и~управление>> Российской академии наук (ФИЦ ИУ РАН);

   Институт проб\-лем информатики ФИЦ ИУ РАН,

AZatsarinny@ipiran.ru}

\footnotetext[2]{Институт проблем информатики Федерального исследовательского

  центра <<Информатика и~управление>> Российской академии наук; 

факультет кибернетики и~информационной безопас\-ности 

Научно-исследовательского ядерного университета МИФИ, vkorolev@ipiran.ru}





\label{st\stat}



                  \thispagestyle{headings}

                  

                  \vspace*{-10pt}

     

     \Abst{Рассмотрено назначение и~характеристика ситуационных центров (СЦ) как 

ин\-фор\-ма\-ци\-он\-но-тех\-но\-ло\-ги\-че\-ских объектов, обеспечивающих реализацию 

ситуационного управ\-ле\-ния. Ситуационный центр представляется сис\-те\-мой. 

Отмечается сложный со\-став и~структура информационных ресурсов, 

обеспечивающих принятие решений. В~качестве фактора проб\-ле\-мы 

информационной без\-опас\-ности (ИБ) СЦ принимается различный 

уровень конфиденциальности обрабатываемой информации. Информационная 

без\-опас\-ность представляется как защита информации. Предлагается подход 

построения сис\-те\-мы защиты информации (СЗИ) для СЦ на основании 

архитектурного подхода. При этом архитектура любой управ\-ля\-емой организации 

представляется как совокупность: биз\-нес-ар\-хи\-тек\-ту\-ра\;+\;архитектура 

информационных технологий (ИТ), 

а~понятие архитектуры сис\-те\-мы соответствует стандарту ГОСТ Р~ИСО/МЭК  

15288-2008. Определены актуальные задачи ИБ

СЦ.}

    

    \KW{ситуационный центр; ситуационное управление; 

информационные ресурсы; ин\-фор\-ма\-ци\-он\-но-ком\-му\-ни\-ка\-ци\-он\-ные технологии; 

конфиденциальность; информационная безопасность; сис\-те\-ма защиты информации; 

архитектурный подход}



\DOI{10.14357\!/\!08696527160109}



\vspace*{-6pt}







\vskip 10pt plus 9pt minus 6pt





      \thispagestyle{headings}

     

    \section{Назначение и~характеристика ситуационных центров}

    

    \vspace*{-2pt}

     

    Ситуационный центр представляется как комплекс, который 

концентрирует информацию о~подконтрольном пространстве от различных 

источников и~обеспечивает ситуационное управление, принятие 

управленческих решений с~широким использованием  

ин\-фор\-ма\-ци\-он\-но-ком\-му\-ни\-ка\-ци\-он\-ных технологий (ИКТ), моделей и~методов 

ситуационного анализа. 

    

    <<Ситуационность>> означает, что действия субъектов, реализация 

процессов и~действий определяются контекстом, в~котором они 

осуществляются. Это понятие лежит в~основе ситуационной теории 

управления, изучающей зависимость эффективности методов управления от 

того, в~каком положении находится объект управления~[1]. 

    

    Ситуационное управление как система становится эффективным, когда 

необходимо, чтобы <<управляющие воздействия субъекта управления 

основывались на анализе вариантов принятия решения с~учетом текущего 

состояния объекта управления, располагаемых вариантов действий 

и~прогноза последствий принимаемых управленческих воздействий>>~[2]. 

    

    Такая необходимость практически постоянно присутствует, когда 

объектами управления являются большие и~сложные объекты, по своей 

природе являющиеся социотехническими сис\-те\-ма\-ми. Ими могут быть 

управленческие, предпринимательские, производственные,  

фи\-нан\-со\-во-опе\-ра\-ци\-он\-ные и~прочие образования (организации, предприятия, 

фирмы, фи\-нан\-со\-во-кре\-дит\-ные учреждения и~т.\,д.), а~так\-же 

муниципальные, региональные и~федеральные образования 

с~соответствующими государственными органами, обеспечивающими 

жизнедеятельность и~управление. 

    

    Объекты управления считаются большими, если они имеют 

распределенную организационную структуру, которая включает в~себя 

компоненты многочисленных видовых групп деятельности с~достаточно 

представительным множеством территориально распределенных однородных 

с точки зрения вида деятельности структурных подразделений. 

    

    Объекты управления считаются сложными, если для реализации их 

целевых функций требуется создание и~функционирование сложного 

и~многопланового инструментария: моделей, процессов и~технологий; 

организационных, тех\-ни\-ко-тех\-но\-ло\-ги\-че\-ских и~инженерных инфраструктур; 

нор\-ма\-тив\-но-пра\-во\-во\-го и~пла\-но\-во-фи\-нан\-со\-во\-го обеспечения; людских 

ресурсов с~широкими группами различной специализации. Живучесть таких 

объектов в~современных условиях зависит от конкурентной борьбы 

и~различного вида противостояний и~поддерживается средствами с~широким 

использованием ИТ~--- признаком 

и~сущностью информационного общества. 

    

    Именно для таких больших и~сложных сис\-тем\-ных объектов существует 

большая вероятность возникновения под влиянием внутреннего и~внешнего 

воздействия совокупности обстоятельств, которые нарушают заданное 

функционирование и~требуют его перевода в~новое состояние, чтобы 

исключить нарушение штатного режима или обеспечить наибольшую 

эффективность функционирования. А~следовательно, становится 

необходимым и~эффективным ситуационное управление.

    

    Если исходить из этого общего представления, СЦ

должен иметь комплексную функциональную направленность обеспечения 

процессов управления и~поддержки принятия решений и~включать в~себя 

следующие компоненты ин\-фор\-ма\-ци\-он\-но-те\-ле\-ком\-му\-ни\-ка\-ци\-он\-ной 

сис\-те\-мы~[3]: 

    \begin{itemize}

\item  комплекс аппаратно-программных средств (средства отображения 

информации, средства хранения информации, средства 

телекоммуникационные, средства защиты информации, автоматизированные 

рабочие места, средства жизнеобеспечения, средства общего 

и~обще\-сис\-тем\-но\-го программного обеспечения (ПО));

    \item комплексы специального ПО, реализующие 

выполнение функциональных задач по назначению СЦ;

    \item обслуживающий персонал, специально подготовленный для 

эффективного применения и~поддержания функционирования СЦ в~заданных 

режимах с~соблюдением мер по ИБ. 

    \end{itemize}

    

    Организационно СЦ состоит из четырех сегментов~[3, 4]:

    \begin{enumerate}[(1)]

    \item руководства органа управления;

    \item мониторинга состояния контролируемого информационного 

пространства;

    \item ситуационного анализа и~поддержки процесса принятия 

управленческих решений, включая ситуационный зал (или кабинет, 

комнату);

    \item администрирования информационного обеспечения,  

ап\-па\-рат\-но-про\-грам\-мных средств и~средств защиты информации. 

    \end{enumerate}

    

    Таким образом, СЦ~--- это базирующийся на  

организационно-техническом комплексе ин\-фор\-ма\-ци\-он\-но-тех\-но\-ло\-ги\-че\-ский 

объект в~сис\-те\-ме управ\-ле\-ния, обеспечивающий реализацию необходимых 

ИКТ, подготовку и~предложение 

ин\-фор\-ма\-ци\-он\-но-ана\-ли\-ти\-че\-ских решений и~поддержку процессов принятия 

решений на компетентном уровне (лица принятия решений~--- ЛПР) 

в~режимах ситуационного управления. На его основе должны быть 

обеспечены мониторинг факторов влияния на развитие происходящих 

процессов, ин\-фор\-ма\-ци\-он\-но-ана\-ли\-ти\-че\-ская поддержка процедур и~процессов, 

позволяющих оперативно анализировать, моделировать, прогнозировать 

сценарии развития ситуаций и~динамично вырабатывать эффективные 

решения.



    \vspace*{-4pt}

    

    \section{Фактор информационной безопасности}

    

        \vspace*{-2pt}

    

    Одним из определяющих факторов выбора регламентов и~обеспечения 

управ\-ле\-ния функционированием СЦ, организации работы каждого из 

перечисленных компонентов, особенностей выбора технических 

и~инженерных решений при создании и~эксплуатации СЦ является 

характеристика конфиденциальности используемых информационных 

ресурсов. Как правило, собственные накапливаемые информационные 

ресурсы СЦ и~привлекаемые внешние информационные ресурсы для 

обеспечения ситуационного управления и~принятия решений независимо от 

вида объекта управления содержат в~той или иной степени информацию, 

требующую контролируемого доступа. В~соответствии с~законодательством 

Российской Федерации это могут быть персональные данные, сведения, 

относящиеся к~определенным видам тайн (государственная, коммерческая, 

банковская, служебная и~др.). И~даже в~общем случае открытые данные, но 

относящиеся к~планируемым инновационным и~инвестиционным проектам, 

к~авторским правам, к~оперативным текущим характеристикам, часто 

в~определенной конкретной ситуации становятся служебной информацией 

и~также требуют контролируемого доступа. Большую роль играет при 

рассмотрении тех или иных ситуаций проблема безопасного агрегирования 

различной информации из собственных информационных ресурсов 

и~внешних источников. 

    

    Исходя из этого, в~современных условиях при использовании для 

принятия решений больших объемов информации, различной по уровню 

конфиденциальности и~признаку владения и~распоряжения, к~СЦ 

предъявляются повышенные требования по обеспечению 

ИБ. При этом под обеспечением ИБ в~данной работе 

понимается решение сис\-тем\-ной задачи защиты информации от 

несанкционированного доступа к~ней, несанкционированной утечки 

информации или ее модификации (нарушения целостности) при обработке 

и~хранении в~ин\-фор\-ма\-ци\-он\-но-ана\-ли\-ти\-че\-ской системе (ИАС) СЦ, передаче 

в~телекоммуникационные системы (ТКС), при реализации технологий 

ситуационного управления и~принятия решений. Хотя это только часть 

сложной и~многоаспектной проблемы обеспечения 

ИБ автоматизированных информационных сис\-тем на объектах 

информатизации~[5].

    

    Следует отметить, что на сегодняшний день такой подход соответствует 

действующим нормативным и~руководящим документам, в~которых 

ИБ при использовании ИТ

определяется как <<защита конфиденциальности, целостности 

и~доступности информации>>~[6].

    

    Поскольку предметом защиты является информация, а~СЦ выступает 

объектом реализации определенных ИТ, то 

сис\-тем\-ная задача защиты информации решается на всем технологическом 

пространстве СЦ, в~его компонентах, при реализации самих процессов 

управления, в~которых участвует СЦ. В~свою очередь, СЦ является 

ма\-те\-ри\-аль\-но-тех\-ни\-че\-ским объектом, и~в~этой части на него 

распространяются положения ГОСТ Р 51275-99~[7], в~котором в~состав 

защищаемого объекта информатизации включаются средства 

жизнеобеспечения (здания, сооружения, помещения, обеспечивающие 

технические средства и~инженерные сис\-те\-мы), которые необходимы для 

установки и~эксплуатации средств и~сис\-тем обработки информации. Особо 

выделяются <<помещения и~объекты, предназначенные для ведения 

конфиденциальных переговоров>>.

    

    С учетом данных предпосылок СЦ формируется как объект защиты, 

и~одним из его сис\-тем\-ных компонентов должен быть компонент 

обеспечения защиты информации.

     

    \vspace*{-4pt}

    

        \section{Общее представление ситуационного центра как системы}



            \vspace*{-2pt}

    

    Ситуационный центр является сложной автоматизированной 

ИАС (АИАС). Необходимо исходить из 

того, что СЦ является частью сис\-те\-мы управ\-ле\-ния объекта (организации как 

объекта), дополняет ее, представляя собой, по существу,  

ин\-фор\-ма\-ци\-он\-но-ана\-ли\-ти\-че\-скую под\-сис\-те\-му или под\-сис\-те\-му поддержки 

принятия решений в~составе сис\-те\-мы управления объекта. При 

декомпозиции СЦ необходимо учитывать сис\-тем\-ную организацию 

и~функциональное наполнение управления объектом как над\-сис\-те\-мы для 

СЦ, а~также учитывать это в~части взаимосвязи всех структурных 

компонентов СЦ с~целью обеспечения эффективного выполнения задач СЦ.

    

    Общая структурно-функциональная схема СЦ как сис\-те\-мы 

представляется в~виде взаимосвязанной совокупности трех базовых 

составляющих: ин\-фор\-ма\-ци\-он\-но-ана\-ли\-ти\-че\-ской,

 ин\-фор\-ма\-ци\-он\-но-тех\-но\-ло\-ги\-че\-ской и~технической~[4, 8]. 

    

    Информационно-аналитическая составляющая включает в~себя 

совокупность методов, методик, моделей, типовых ситуаций, возможных 

реакций на угрозы, сценариев принятия решений, ориентированных на 

решение задач ситуационного анализа в~рамках заданной предметной 

области СЦ. Практически это сервисные средства ситуационного 

управления, которые должны обеспечивать деятельность  

экс\-перт\-но-ана\-ли\-ти\-че\-ско\-го и~оперативного состава органов управления 

(экспертно-аналитическая и~дежурная службы), отвечающих за разработку 

и~выдачу ин\-фор\-ма\-ци\-он\-но-ана\-ли\-ти\-че\-ских продуктов для ЛПР 

и~за подготовку и~сопровождение процессов ситуационного 

управления. 

    

    Инструментально информационно-аналитическая составляющая 

представляет собой приложения прикладного ПО и~регламенты работы 

с~ним, предназначенные для решения функциональных задач СЦ. Это 

средства ситуационного анализа, моделирования ситуаций, сис\-те\-мы задания, 

распознавания и~извлечения из информационных ресурсов необходимой 

информации, формирования информационных продуктов визуализации 

и~решения ин\-фор\-ма\-ци\-он\-но-рас\-чет\-ных задач, выбора и~настройки сценариев 

и~регламентов технологического процесса ситуационного управления. 

    

    Информационно-технологическая составляющая~--- это исполнительная 

ин\-фор\-ма\-ци\-он\-но-тех\-но\-ло\-ги\-че\-ская платформа СЦ. Она должна обеспечить 

необходимыми информационными ресурсами и~реализовывать базовые 

технологии (технологические сервисы, процедуры) функционирования СЦ 

как части сис\-те\-мы управ\-ле\-ния объекта, использовать в~соответствующих 

сценариях ситуационного управления ин\-фор\-ма\-ци\-он\-но-ана\-ли\-ти\-че\-ские 

продукты, настройки и~регламенты, подготавливаемые для этих сценариев. 

    

    Одной из базовых функций ин\-фор\-ма\-ци\-он\-но-тех\-но\-ло\-ги\-че\-ской 

составляющей является формирование и~ведение информационных ресурсов 

СЦ. Это, прежде всего, сбор из различных источников данных 

(формализованной и~неформализованной информации), обработка, анализ, 

обобщение и~отображение данных в~базах данных применительно к~циклу 

<<со\-бы\-тие--си\-ту\-а\-ция--угро\-за>>, интеграция и~предоставление 

данных для решения функциональных задач, а~также управление 

внутренними и~внешними информационными потоками. Ведение 

информационных ресурсов СЦ требует ведения единого лингвистического 

обеспечения.

    

    Инструментально информационно-тех\-но\-ло\-ги\-че\-ская со\-став\-ля\-ющая 

содержит комплексы специального ПО, а~также 

общее ПО, выбор которого обуслов\-ли\-ва\-ет\-ся 

спецификой функционирования СЦ и~требованиями по защите информации. 

Очевидно, что эта со\-став\-ля\-ющая определяет целевое предназначение 

и~эффективность сис\-те\-мы СЦ в~целом. 



     \begin{figure}[b] %fig1

     \vspace*{-6pt}

 \begin{center}

 \mbox{%

 \epsfxsize=127.088mm

 \epsfbox{zac-1.eps}

 }

\end{center}

 \vspace*{-11pt}

     \Caption{Обобщенная структурно-функциональная схема сис\-те\-мы СЦ}

     \end{figure}

    

    Между информационно-аналитической  

и~ин\-фор\-ма\-ци\-он\-но-тех\-но\-ло\-ги\-че\-ской составляющими существует 

непрерывное технологическое и~информационное взаимодействие, обмен 

ресурсами, и~в~этой части существует постоянная потребность уточнения их 

корреляции, которая, как правило, определяется и~осуществляется в~рамках 

ИАС СЦ. 





    

    Техническая составляющая является ап\-па\-рат\-но-про\-граммной 

платформой СЦ, обеспечивающей вычислительную 

и~телекоммуникационную технологическую среду для реализации задач СЦ 

и~его функционирования в~целом.

    

    Инструментально техническая составляющая включает в~себя 

вычислительный комплекс ИАС, ТКС с~оборудованием визуализации 

информации и~видеоконференцсвязи, оборудование информационного 

взаимодействия с~внешними ИАС и~базами данных, в~том числе в~режиме 

облачного взаимодействия, их обще\-сис\-тем\-ное ПО

и~обще\-сис\-тем\-ные программные сервисы. К~общесистемным программным 

сервисам следует отнести и~сервисы управления компонентами и~процессами 

ап\-па\-рат\-но-про\-грам\-мной платформы~СЦ.

    

    Согласно ГОСТ Р 51275-99~[7] в~состав технической составляющей СЦ 

включаются средства жизнеобеспечения (здания, сооружения, помещения, 

обеспечивающие технические средства и~инженерные сис\-те\-мы), которые 

необходимы для установки и~эксплуатации средств и~сис\-тем  

ап\-па\-рат\-но-про\-граммной платформы СЦ, в~том числе помещения, 

предназначенные для ведения конфиденциальных переговоров при 

функционировании СЦ.

    

    Информационно-технологическая и~техническая составляющие 

формируют базовую ин\-фор\-ма\-ци\-он\-но-те\-ле\-ком\-му\-ни\-ка\-ци\-он\-ную 

инфраструктуру СЦ. Служба эксплуатации обеспечивает поддержание 

средств базовой информационно-те\-ле\-ком\-му\-ни\-ка\-ци\-он\-ной инфраструктуры 

в~требуемом состоянии и~их штатное применение. 

    

    Информационная безопасность СЦ обеспечивается СЗИ, 

    которая комплексно решает задачи защиты для СЦ в~целом. 

    

    Обобщенная структурно-функциональная схема сис\-те\-мы СЦ 

представлена на рис.~1. 

    

    \vspace*{-4pt}

    

    \section{Архитектурный подход к~проектированию ситуационного центра 

и~обеспечению информационной безопасности} 



    \vspace*{-2pt}

    

    Будучи частью системы управления организации в~целом, СЦ 

представляет собой ин\-фор\-ма\-ци\-он\-но-тех\-но\-ло\-ги\-че\-ский объект. Его 

проектирование, с~одной стороны, тесно связано с~биз\-нес-про\-цес\-са\-ми, 

деловыми и~служебными процессами деятельности организации, с~другой 

стороны, СЦ проектируется как сис\-тем\-ный компонент ИТ-ар\-хи\-тек\-ту\-ры 

организации. Такая схема построения типична в~современных условиях 

глубокого проникновения ИКТ в~биз\-нес-про\-цес\-сы и~управление и~требует 

архитектурного подхода при проектировании,\linebreak когда под архитектурой 

организации понимается совокупность:  

биз\-нес-ар\-хи\-тек\-ту\-ра\;+\;ар\-хи\-тек\-ту\-ра ИТ~[9]. 

    %

    При этом будем придерживаться понятия, соответствующего стандарту 

ГОСТ Р~ИСО/МЭК 15288-2008: архитектура сис\-те-\linebreak мы~--- принципиальная 

организация сис\-те\-мы, воплощенная в~ее элементах, их взаимоотношениях 

друг с~другом и~со средой, а~так\-же принципы, направляющие ее 

проектирование и~эволюцию. 



В~то же время следует учитывать и~свойство 

рекурсивности в~понятии архитектуры сис\-те\-мы, часто используемое 

в~области создания сложных программных продуктов: архитектура сис\-те\-мы 

состоит из нескольких компонентов, внешних свойств и~интерфейсов, связей 

и накладываемых ограничений, а~также архитектуры этих внутренних 

компонентов.

    

    Архитектурный подход позволяет проектировать обеспечение ИБ СЦ 

как частную архитектуру сис\-тем\-ной архитектуры организации, и~ее 

построение должно соответствовать или, по крайней мере, не противоречить 

архитектурным, структурным и~техническим решениям по автоматизации 

и~информатизации организации~[10].

    

    \vspace*{-4pt}

    

        \section{Система защиты информации ситуационного центра}

        

            \vspace*{-2pt}

     

    Реализация мер организационного, технологического и~технического 

характера по защите информации в~СЦ возлагается на 

СЗИ, которая при проектировании СЦ как сис\-те\-мы 

является ее \textit{обеспечивающей под\-сис\-те\-мой}~[4]. 

    

    В рамках архитектурного подхода к~созданию СЦ базовым положением 

проектирования СЗИ выступает следующее требование: структурные, 

технические, технологические и~организационные решения по защите 

информации должны быть неразрывно связаны с~проектными решениями по 

СЦ и~обеспечивающей его функционирование инфраструктуре. Это 

требование служит предпосылкой для сис\-тем\-ной реализации ИБ 

и~комплексного использования всех известных и~соответствующих 

нормативам методов и~средств защиты информации, т.\,е.\ создания 

комплексной СЗИ СЦ, которая осуществляет сис\-тем\-ную интеграцию 

разнородных задач и~средств для достижения единой цели~--- обеспечения 

заданного уровня защиты. 

    

    Средства защиты, входящие в~СЗИ СЦ, должны регулярно выполнять 

свои функции, не конфликтуя между собой и~с~другими компонентами, 

обеспечивающими функционирование СЦ. Они могут быть реализованы 

различными способами и~методами в~части моделей, алгоритмов 

и~инструментария, могут быть программными или аппаратными средствами, 

организационными мерами и~мероприятиями, а~также представлять собой 

комплексные технические и~организационные решения в~виде отдельных 

функциональных изделий.

    

    Системное решение обеспечения ИБ в~виде СЗИ предполагает при 

функционировании СЦ обеспечение защиты информации~[11]:

    \begin{itemize}

    \item при всех видах информационной деятельности на объекте 

управления, в~которых используется СЦ;

    \item во всех структурных компонентах СЦ;

    \item при всех режимах функционирования СЦ;

    \item на всех этапах жизненного цикла объекта управления 

и,~соответственно, СЦ;

    \item  с~учетом информационного взаимодействия с~внешней средой.

    \end{itemize}

    

    В общем случае СЗИ СЦ разрабатывается как под\-сис\-те\-ма 

автоматизированной сис\-те\-мы, так как ранее было отмечено, что СЦ~--- это 

сложная АИАС. 

Типовые подходы к~проектированию сис\-тем защиты подобного назначения 

в~настоящее время практически сложились~[5], однако имеют место свои 

особенности и~новые проблемы по ИБ, которые 

присущи СЦ как ин\-фор\-ма\-ци\-он\-но-тех\-но\-ло\-ги\-че\-ско\-му и~функциональному 

компоненту в~сис\-те\-ме управ\-ления. 

    

    По системному представлению СЗИ СЦ является автоматизированной 

ор\-га\-ни\-за\-ци\-он\-но-тех\-ни\-че\-ской сис\-те\-мой управ\-ле\-ния ИБ и~включает в~себя 

функциональную и~обеспечивающую части. 

    

    \textit{Функциональная часть СЗИ СЦ}~--- это совокупность 

функциональных задач по нейтрализации угроз информации, которые могут 

иметь различные источники и~быть различной природы. 

\textit{Функциональные задачи} ставятся и~решаются в~отношении 

актуальных для СЦ угроз. При этом обеспечивается:

    \begin{itemize}

    \item рассмотрение всех угроз информации, влияющих на состояние 

защиты информации в~СЦ; 

    \item выделение из них актуальных путем анализа порождаемых ими 

рисков и~последующего возможного ущерба, если риск становится реальным 

событием;

    \item ранжирование актуальных угроз;

    \item введение в~классификационную схему угроз информации всех 

выделенных актуальных угроз, учитываемых при создании СЗИ СЦ.

    \end{itemize}

    

    Функциональные задачи защиты образуют \textit{функциональные 

под\-сис\-те\-мы СЗИ СЦ}, объединяющие задачи по признакам их решения 

в~определенных ИТ, компонентах СЦ или по 

принадлежности группы функциональных задач к~выделенному направлению 

защиты информации.

    

    Примерный состав функциональных под\-сис\-тем СЦ, который, 

безусловно, должен быть конкретизирован и~обоснован для каждого 

проектируемого СЦ, может быть следующим, а~именно включать 

под\-сис\-темы:

    \begin{itemize}

    \item  защиты информации в~ТКС;

    \item защиты от несанкционированного доступа (НСД) к~информации 

и~ресурсам в~ИАС, в~том числе управления разграничением доступа 

к~информации, ресурсам, режимам функционирования ИАС и~режимам 

ситуационного управ\-ления;

    \item обеспечения ИБ операционного 

ситуационного помещения (за\-ла\!/\!ком\-на\-ты\!/\!ка\-би\-не\-та), в~том 

числе защиты информации визуализации и~видеоконференцсвязи;

    \item контроля и~разграничения доступа при внешнем и~удаленном 

информационном взаимодействии, в~том числе в~регламенте облачного 

взаимодействия с~внешними базами данных;

    \item секретного и~конфиденциального электронного документооборота 

(ЭДО);

    \item обеспечения целостности информации и~цифровой подписи, в~том 

числе создания инфраструктуры открытых ключей для СЦ 

и~взаимодействующих с~ним субъектов и~объектов;

    \item мониторинга и~противодействия сетевым компьютерным атакам, 

в~том числе защиты от вирусов;

    \item создания и~поддержания доверенной про\-граммно-тех\-ни\-че\-ской 

среды, в~том числе предотвращения скрытного внедрения в~программные 

и~технические средства;

    \item защиты речевой информации в~процессах ситуационного 

управления и~принятия решений;

    \item защиты информации от утечки по техническим каналам, в~том 

числе по техническим проводным каналам инженерных обеспечивающих 

сис\-тем;

    \item контроля защищаемого контура размещения служб и~средств СЦ 

и~другие возможные под\-сис\-темы.

    \end{itemize}

    

    Важным системным требованием является требование обеспечения 

ИБ на всем технологическом пространстве СЦ.

    %

    При этом функция защиты определенного вида и~назначения, 

относящаяся к~выделенной функциональной под\-сис\-те\-ме СЗИ, может 

выполняться в~различных компонентах СЦ, технологических процессах 

и~точках про\-граммно-тех\-ни\-че\-ской среды (например, функция 

аутентификации).

    

    Другим распространенным вариантом композиционного соотношения 

функций защиты является вариант, при котором одна и~та же функция 

защиты должна выполняться в~различных функциональных под\-сис\-те\-мах 

СЗИ, оставаясь в~то же время функцией выделенной функциональной 

под\-сис\-те\-мы СЗИ (например, мониторинг и~противодействие сетевым 

атакам).

    

    Чтобы избежать такого структурного противоречия, при построении 

СЗИ СЦ продуктивно включать в~функциональную часть 

\textit{функциональные комплексы защиты информации}, которые 

становятся самостоятельными компонентами СЗИ СЦ. При этом интеграция 

однородных функциональных задач в~рамках под\-сис\-те\-мы СЗИ образует 

функциональный комплекс I~вида, а~интеграция однородных 

функциональных задач, реализуемых в~отношении одной и~той же функции 

назначения в~различных под\-сис\-те\-мах СЗИ, образует функциональный 

комплекс~II~вида.

    

    Выделение функциональных комплексов при создании СЗИ 

и~последующей ее эксплуатации~--- это очевидный путь к~типовым 

решениям по применению или созданию средств защиты информации 

с~учетом включенных в~эти средства возможностей настройки на параметры 

политики безопасности и~информационного взаимодействия в~части сбора 

служебной контролирующей информации для реализации управления. 

    

    Схема формирования функциональных комплексов показана на рис.~2.

     



     

    \textit{Обеспечивающая часть СЗИ СЦ}~--- совокупность мер и~средств, 

технических, организационных и~правовых решений, обеспечивающих 

реализацию функциональных задач.

    

    К обеспечивающей части предъявляются обще\-сис\-тем\-ные 

требования~[12]:

    \begin{itemize}

    \item должны быть разработаны и~доведены до уровня регулярного 

использования все необходимые механизмы гарантированного обеспечения 

требуемого уровня защиты информации; 

    \item механизмы требуемого уровня защиты информации должны 

существовать в~практической реализации в~виде средств и~технологий; 

    \item необходимо располагать сертифицированными средствами (или 

отвечающими требованиям обработки конфиденциальной информации 

соответствующего уровня) реализации всех необходимых технических 

решений по защите информации; 

    \item должны быть разработаны способы, решения и~условия 

оптимальной организации и~обеспечения проведения всех мер и~мероприятий 

по защите в~процессе обработки информации;

    \item должно быть создано правовое поле для реализации функций 

защиты, механизмы и~средства реализации правовых положений 

и~требований при функционировании средств защиты и~СЗИ СЦ в~целом.

    \end{itemize}

    

         \begin{figure} %fig2

     \vspace*{1pt}

 \begin{center}

 \mbox{%

 \epsfxsize=114.641mm

 \epsfbox{zac-2.eps}

 }

\end{center}

 \vspace*{-11pt}

    \Caption{Схема формирования функциональных комплексов СЗИ СЦ}

    \vspace*{-6pt}

     \end{figure}

    

    Системная интеграция функциональных задач и~средств их реализации 

обеспечивается, наряду с~\textit{технологическим обеспечением СЗИ, 

управлением} как сис\-тем\-ным компонентом СЗИ. Основанием для принятия 

решений при управ\-ле\-нии служат результаты мониторинга со\-сто\-яния 

ИБ. Как правило, в~СЗИ СЦ выделяется 

\textit{управ\-ле\-ние процессами защиты} и~\textit{менеджмент обеспечения 

ИБ в~СЦ} в~целом~[10].

    

    Технологическая интеграция и~управление должны обеспечить 

достижение следующих основных целей~[4]:

    \begin{itemize}

    \item предотвращение утечки информации по техническим каналам, 

в~том числе речевой информации, а~также утечки за счет возможно 

внедренных специальных устройств негласного получения информации; 

    \item предотвращение несанкционированного уничтожения, искажения, 

копирования и~блокирования информации за счет НСД к~ресурсам и~сетевых атак; 

    \item соблюдение правового режима использования информационных 

ресурсов и~программ обработки информации;

    \item обеспечение целостности информации при ее получении, 

обработке и~хранении в~СЦ; 

    \item сохранение управляемости процессами обработки, хранения 

и~использования информации в~условиях несанкционированных воздействий 

на защищаемую информацию и~среду ее обработки;

    \item  обеспечение защищенного обмена информацией между объектами 

СЦ, в~том числе удаленного и~внешнего информационного обмена. 

    \end{itemize}

    

    \begin{figure}[b] %fig3

\vspace*{1pt}

 \begin{center}

 \mbox{%

 \epsfxsize=126.449mm

 \epsfbox{zac-3.eps}

 }

\end{center}

 \vspace*{-9pt}

     \Caption{Схема структурных отношений между компонентами СЗИ СЦ}

     \end{figure} 

    

    Проектирование и~последующая эксплуатация СЗИ СЦ должна 

базироваться на концептуальных нормативных документах: \textit{Модели 

угроз информационной безопасности}, \textit{Модели вероятного 

нарушителя} и~\textit{Политики информационной безопасности СЦ}. 

Модели угроз и~вероятного нарушителя являются исходными 

концептуальными документами, подготавливаемыми в~основном по 

результатам предпроектного обследования. Они позволяют определить 

актуальные угрозы, уязвимости, должны ответить на вопросы~--- что 

защищать и~от кого защищать, конкретизировать постановку задач по защите 

информации и~тем самым обеспечить минимизацию затрат на ИБ.

    

    Политика ИБ в~своей основе рассчитана на 

время жизни СЦ. В~виде нормативного сис\-тем\-но\-го документа она 

предъявляется на этапе при\-емо-сда\-точ\-ных испытаний СЦ с~возможной 

последующей корректировкой в~ходе эксплуатации. 

    

    Схема структурных отношений между компонентами СЗИ СЦ 

представлена на рис.~3.

     

    \vspace*{-4pt}

     

    \section{Актуальные задачи информационной безопасности ситуационного центра}

    

        \vspace*{-2pt}

     

    Ситуационные центры являются объектами информационной 

деятельности с~отчетливо выраженной спецификой. Они должны 

обеспечивать безопасное агрегирование получаемых из различных 

источников больших объемов различной информации и~ее  

ин\-фор\-ма\-ци\-он\-но-ана\-ли\-ти\-че\-скую обработку средствами собственной 

инфраструктуры. В~СЦ стекается информация, разнородная не только по 

своей сути и~форме (документы, потоковые данные, телеметрия и~пр.), но 

и~по уровню конфиденциальности~--- <<грифу>>. В~связи с~этим общие 

требования по ИБ к~СЦ формируются на основе 

требований к~защищаемым ресурсам с~наивысшим обрабатываемым 

<<грифом>>.

    

    В связи со спецификой СЦ, их усложнением и~бурным развитием 

в~настоящее время появляются новые актуальные задачи по обеспечению 

ИБ СЦ~[13--15]. Практическое решение этих 

задач все более подтверждает тот факт, что функции защиты информации 

становятся неотъемлемой частью функциональных сервисов обработки 

информации и~архитектурных решений при по\-стро\-ении сис\-тем\-ных 

инфраструктурных компонентов. Это, безусловно, влияет на проектные 

решения СЗИ СЦ.

    

    К таким актуальным задачам можно отнести:

    \begin{itemize}

    \item  \textit{сегментирование защищаемых ресурсов СЦ} 

(информационных, ин\-фор\-ма\-ци\-он\-но-тех\-но\-ло\-ги\-че\-ской инфраструктуры 

и~про\-граммно-тех\-ни\-че\-ской среды)~--- выделение контуров безопасности по 

<<грифу>>. Это решение считается одним из доверенных механизмов 

защиты для накопительных информационных сис\-тем с~большими объемами 

разнородной поступающей информацией;

    \item  \textit{создание защищенной шины информационного 

взаимодействия}, способной объединить разнородные первичные ресурсы 

в~случае их выраженной территориальной распределенности. Применение 

терминального доступа или использование кли\-ент-сер\-вер\-ных технологий 

в~этом случае неэффективно, поскольку необходима реализация высоких 

требований по пропускной способности и~качеству каналов связи, 

образующих единую сеть;

    \item \textit{построение СЦ как частного облака}, когда возникает 

проблема виртуального пространства взаимодействия. Это приводит 

к~усложнению СЗИ СЦ с~учетом специфических угроз, присущих 

виртуальной среде исполнения;

    \item \textit{организацию эргономичного рабочего места} 

с~минимизацией многооконного интерфейса, не вполне удобного 

и~понятного для ЛПР. Интеграция разнокатегорированных сис\-тем~--- 

<<генераторов>>~--- в~рамках <<одного окна>> является сложнейшей 

проблемой не только с~точки зрения обеспечения ИБ, но и~с~инженерной 

точки зрения;

    \item \textit{безопасное хранение в~ИАС СЦ поколений разнородной, 

в~том числе графической и~потоковой, информации}, что предъявляет жесткие 

требования по ИБ к~сис\-те\-мам управ\-ле\-ния базами данных (СУБД). Это требует применения сертифицированных 

СУБД или использования наложенных средств контроля доступа. Последнее 

решение существенно сужает возможности и~самой СУБД, и~средств 

углубленного анализа;

    \item \textit{реализацию режима чрезвычайной ситуации}, когда все 

внимание СЦ на\-прав\-ля\-ется на формирование новостного информационного 

потока исключительно вокруг интересующего события. В~этом случае 

уровень конфиденциальности СЦ может быть повышен до максимального 

уровня. Возможен разрыв стыков с~открытыми контурами СЦ. Ситуационный центр может 

стать автономным объектом, взаимодействие с~которым осуществляется либо 

по доверенным каналам связи, либо посредством подкачки данных 

с~отчуждаемых носителей, физически доставляемых в~СЦ;

    \item \textit{формирование электронного указания} в~виде приказа или 

распоряжения как способ управляющего воздействия. Это потребует 

электронного документооборота, обеспечивающего возможность 

статуса правового подтверждения, что, в~свою очередь, потребует 

сопряжения СЦ с~удостоверяющим центром и~создания 

инфраструктуры открытых ключей для всего сетевого сообщества СЦ.

    \end{itemize}

    

    Следует отметить, что перечисленные и~ряд других актуальных задач 

обеспечения ИБ СЦ, особенно связанных 

с~планируемым созданием в~России сети распределенных 

СЦ~[8, 16], требуют своего достаточно быстрого и~качественного 

решения, чтобы поддержать большую востребовательность СЦ для решения 

задач управления в~современном быстро меняющемся обществе.



    \vspace*{-4pt}

      

{\small\frenchspacing

{%\baselineskip=10.8pt

\addcontentsline{toc}{section}{Литература}

\begin{thebibliography}{99}



    \vspace*{-2pt}

    

\bibitem{1-kor}

\Au{Маршев В.\,И.} История управленческой мысли.~--- М.: Инфра-М, 2005. 731~с.

\bibitem{2-kor}

Словарь терминов МЧС.~--- М.: Федеральное государственное унитарное авиационное 

предприятие МЧС России EdwART, 2010. 

\bibitem{3-kor}

\Au{Зацаринный А.\,А., Сучков А.\,П.} Системы ситуационных центров специального 

назначения. Основные определения, понятия и~подходы к~созданию~/\!\!/ Меж\-отрас\-ле\-вая 

информационная служба, 2015. №\,4. С.~31--41.

\bibitem{4-kor}

\Au{Зацаринный А.\,А.} Организационные и~сис\-те\-мо\-тех\-ни\-че\-ские подходы к~построению 

современных ситуационных центров~/\!\!/ Методы построения и~технологии 

функционирования ситуационных центров: Сб. науч.-технич. статей~/ Под ред.  

А.\,А.~Зацаринного.~--- М.: ИПИ РАН, 2011. С.~10--25. 

\bibitem{5-kor}

\Au{Королёв В.\,И.} Методология построения комплексной защиты информации на 

объектах информатизации~/\!\!/ Системы высокой доступности, 2009. Т.~5. №\,4. С.~4--24.

\bibitem{6-kor}

ГОСТ Р ИСО/МЭК 27002-2012. Информационная технология. Методы и~средства 

обеспечения безопасности. Свод норм и~правил менеджмента информационной 

безопасности. 

\bibitem{7-kor}

ГОСТ Р 51275-99. Защита информации. Объект информатизации. Факторы, 

воздействующие на информацию. 

\bibitem{8-kor}

\Au{Зацаринный А.\,А.} Ситуационные центры как основа информационно-аналитической 

поддержки принятия решений в~органах государственной власти~/\!\!/ Аналитика 

развития и~безопасности страны: реалии и~перспективы: Сб. мат-лов 1-й Всеросс. 

конф.~--- М.: Столица, 2014. C.~277--295. 

\bibitem{9-kor}

Подход системной инженерии к~управлению жизненным циклом. Понятийный минимум. 

PraxOS. 1.0, 2008. {\sf http://techinvestlab.ru/files/495344/se\_ls\_\linebreak minimum\_praxos\_1.doc}. 

\bibitem{10-kor}

\Au{Королёв В.\,И.} Концептуальные аспекты построения информационной безопасности 

корпораций~/\!\!/ Системы высокой доступности, 2015. Т.~11. №\,2. С.~50--64.

\bibitem{11-kor}

\Au{Герасименко В.\,А.} Защита информации в~автоматизированных сис\-те\-мах обработки 

данных.~--- М.: Энергоатомиздат, 1994.  Кн.~1. 400~с. 

\bibitem{12-kor}

\Au{Королёв В.\,И., Сухотин И.\,Н., Королёв~А.\,В.} Интегрированные сис\-те\-мы 

без\-опас\-ности и~влияние информационных рисков на деятельность организации~/\!\!/ 

Технологии безопасности: Официальный отчет Х Международного форума: Сб. мат-лов~/ 

Под ред. И.\,К.~Филоненко, Н.\,В.~Александровой.~--- М.: ПРОЭКСПО, 2005.  

С.~284--291. 

\bibitem{13-kor}

\Au{Андреев В.} Защита информации в~ситуационном центре~--- ключевые аспекты: 

Типовые и~специальные задачи~/\!\!/ Connect, 2012. №\,5. {\sf 

http://www.ситцентр.рф/ pressa/andreev.pdf}. 



\bibitem{15-kor}

\Au{Шарай В.\,А., Малашихин А.\,К.} Логико-ве\-ро\-ят\-ност\-ные модели угроз 

информационной безопасности ситуационных центров~/\!\!/ Научные труды КубГТУ, 

2014. №\,6. {\sf http://ntk.kubstu.ru/file/306}. 



\bibitem{14-kor}

\Au{Никифоров Д.} Об информационной безопасности будущей сис\-те\-мы ситуационных 

центров в~России~/\!\!/ Connect WIT, 2015. №\,7-8. {\sf  

http://www.connect-wit. ru/category/magazine/connect-wit-2015-7-8}.



\bibitem{16-kor}

Стратегия национальной безопасности Российской Федерации. Утверждена Указом 

Президента РФ от 31.12.2015 №\,683. Ст.~112.

      \end{thebibliography}

} }





\vspace*{-3pt}



\hfill{\small\textit{Поступила в~редакцию 09.02.16}}





\newpage



%\vspace*{9pt}



%\hrule



%\vspace*{2pt}



%\hrule



\vspace*{-24pt}







\def\tit{INFORMATION SECURITY SITUATION CENTER}



\def\titkol{Information security situation center}





\def\aut{A.\,A.~Zatsarinny$^1$ and~V.\,I.~Korolev$^{1,2}$}

\def\autkol{A.\,A.~Zatsarinny and~V.\,I.~Korolev}







\titel{\tit}{\aut}{\autkol}{\titkol}



\vspace*{-9pt}





\noindent

$^1$Institute of Informatics Problems, Federal Research Center 

``Computer Science\linebreak

$\hphantom{^1}$and Control'' of the Russian Academy of Sciences, 44-2 

Vavilova Str., Moscow\linebreak

$\hphantom{^1}$119333, Russian Federation



\noindent

$^2$Department ``Cybernetics and Information Security,'' National Research 

Nuclear\linebreak

$\hphantom{^1}$University ``MEPhI,'' 31~Kashirskoye Sh., Moscow 115409, 

Russian Federation





\def\leftfootline{\small{\textbf{\thepage}

\hfill Sistemy i Sredstva Informatiki~---

Systems and Means of Informatics 2016 vol~26 no\,1}

}%

 \def\rightfootline{\small{Sistemy i Sredstva Informatiki~---

 Systems and Means of Informatics   2016   vol~26   no\,1

\hfill \textbf{\thepage}}}      

     

     



\Abste{The paper considers the purpose and characteristics of  

situational centers as information technology facilities that ensure the 

implementation of situational management. A~situational center is a~system. The 

paper indicates complicated composition and structure of information resources 

used to make decisions. The variety of levels of confidentiality related to 

the processed information is a~factor of situational centers information security. 

Information security is performed as information protection. The paper presents 

an approach to building the systems of information security and protection for situational 

centers based on the architectural approach. Herewith,  this architecture of  any run 

organization is of a set of business architecture\;+\;IT 

architecture and the concept of system architecture complies with ISO/IEC~15288-2008. 

The paper formulates the actual  problems of information security situational centers.}



\KWE{situational center; situational management; informational resources; 

information and communication technologies; confidentiality; 

information security; system of information protection; architectural approach}









\DOI{10.14357\!/\!08696527160109} %



%\vspace*{-24pt}



\Ack

\noindent

The work was supported by the Russian Foundation for Basic Research 

(project 15-29-07981~ifi-m).





\renewcommand{\bibname}{\protect\rmfamily References}

%\renewcommand{\bibname}{\large\protect\rm References}



{\small\frenchspacing

{%\baselineskip=10.8pt

\addcontentsline{toc}{section}{References}

\begin{thebibliography}{99} 

     

\bibitem{1-kor-1}

\Aue{Marshev, V.\,I.} 2005. \textit{Istoriya upravlencheskoy mysli}. [A~story of 

administrative thought]. Moscow: Infra-M. 731~p.

\bibitem{2-kor-1}

\textit{Slovar' terminov MChS}

[Dictionary of terms of the Ministry of Emergency Situations]. 2010.  Federal 

State Unitary Air Enterprise of Emercom of Russia of EdwART.

\bibitem{3-kor-1}

\Aue{Zatsarinny, A.\,A., and A.\,P.~Suchkov}. 2015. Sistemy situatsionnykh tsentrov 

spe\-tsi\-al'\-no\-go naznacheniya. Osnovnye opredeleniya, ponyatiya i~podkhody k~sozdaniyu 

[Systems of the situational centers of a special purpose. Main definitions, concepts and 

approaches to creation]. \textit{Mezhotraslevaya Informatsionnaya Sluzhba} [Interindustry 

Information Service]  4:31--41. 

\bibitem{4-kor-1}

\Aue{Zatsarinny, A.\,A.} 2011. Organizatsionnye i~sistemotekhnicheskie podkhody 

k~postroeniyu sovremennykh situatsionnykh tsentrov [Organizational and sistemotekhnical 

approaches to creation of the modern situational centers]. \textit{Sb. nauch.-tekhnich. statey 

``Metody postroeniya i~tekhnologii funktsionirovaniya situatsionnykh tsentrov''} [Methods of 

construction and technology of functioning of the situational centers, Collection of Scientific and 

Technical Articles]. Ed.\  A.\,A.~Zatsarinniy. Moscow: IPI RAN. 10--25.

\bibitem{5-kor-1}

\Aue{Korolev, V.\,I.} 2009. Metodologiya postroeniya kompleksnoy zashchity 

informatsii na 

ob'ektakh informatizatsii 

[Methodology of creation of complex information security on objects 

of informatization]. \textit{Sistemy Vysokoy Dostupnosti} 

[Systems of High Availability] 4(5):4--24. 

\bibitem{6-kor-1}

GOST R ISO\,/\,IEC 27002-2012. Informatsionnaya tekhnologiya. Metody i~sredstva 

obespecheniya bez\-opas\-nosti. Svod norm i~pravil menedzhmenta informatsionnoy 

bez\-opas\-nosti 

[Information technology. Methods and means of ensuring of safety. Arch of norms and rules of 

management of information security]. 

\bibitem{7-kor-1}

GOST R 51275-99. Zashchita informatsii. Ob'ekt informatizatsii. Faktory, vozdeystvuyushchie 

na informatsiyu [Information security. Object of informatization. The factors influencing 

information].

\bibitem{8-kor-1}

\Aue{Zatsarinny, A.\,A.} 2014. Situatsionnye tsentry kak osnova informatsionno-analiticheskoy 

podderzhki prinyatiya resheniy v organakh gosudarstvennoy vlasti [Situational centers as a basis 

of information and analytical support of decision-making in Government]. 

\textit{Sb. mat-lov 

1-y Vseross. Konf. ``Analitika razvitiya i~bezopasnosti strany: 

Realii i~perspektivy''} 

[Collection of materials of the 1st All-Russian Conference ``Analytics of Development and 

Safety of the Country: Realities and Prospects'']. Moscow: Stolitsa. 277--295. 



\bibitem{9-kor-1}

PraxO.S. 1.0, 2008. Podkhod sistemnoy inzhenerii k~upravleniyu zhiznennym tsiklom. 

Ponyatiynyy minimum [Approach of system engineering to management of life cycle. 

Conceptual minimum]. Available at: {\sf 

http://techinvestlab.ru/files/ 495344/se\_ls\_minimum\_praxos\_1.doc} (accessed February~9, 

2016). 

\bibitem{10-kor-1}

\Aue{Korolev, V.\,I.} 2015. Kontseptual'nye aspekty postroeniya informatsionnoy 

bez\-opas\-nosti 

korporatsiy [Conceptual aspects of creation of information security of corporations]. 

\textit{Sistemy Vysokoy Dostupnosti} [Systems of High Availability] 2(11):50--64. 

\bibitem{11-kor-1}

\Aue{Gerasimenko, V.\,A.} 1994. \textit{Zashchita informatsii v~avtomatizirovannykh sistemakh 

obrabotki dannykh}  [Information security in the automated systems of data processing]. 

Moscow: Energoatomizdat. Vol.~1. 400~p.

\bibitem{12-kor-1}

\Aue{Korolev, V.\,I., I.\,N. Sukhotin, and A.\,V.~Korolev}. 2005. Integrirovannye sistemy 

bezopasnosti i~vliyanie informatsionnykh riskov na deyatel'nost' organizatsii [The integrated 

security systems and influence of information risks on activity of the organization]. 

\textit{Ofitsial'nyy otchyot X~Mezhdunar. Foruma ``Tekhnologii Bezopasnosti.'' Sb. mat-lov}.

[The Official Report of the 10th Forum (International) ``Technologies of Safety.'' The collection of 

Materials]. Eds.\ I.\,K.~Filonenko and N.\,V.~Aleksandrova. Moscow: PROEKSPO. 

284--291.

\bibitem{13-kor-1}

\Aue{Andreyev, V.} 2012. Zashchita informatsii v~situatsionnom tsentre~--- klyuchevye aspekty. 

Tipovye i~spetsial'nye zadachi [Information security in the situational center~--- key aspects. 

Standard and special tasks]. \textit{Connect} 5. Available at: {\sf 

http:// www.ситцентр.рф/pressa/andreev.pdf} (accessed February~9, 2016). 



\bibitem{15-kor-1}

\Aue{Sharay, V.\,A., and A.\,K.~Malashikhin}. 2014. Logiko-veroyatnostnye modeli ugroz 

informatsionnoy bezopasnosti situatsionnykh tsentrov [Logic and probabilistic model of threats 

of information security of the situational centers]. \textit{Nauchnye trudy KubGTU} [Scientific 

works of KubGTU]  6. Available at: {\sf http://ntk.kubstu.ru/file/306} (accessed February~9, 

2016).



\bibitem{14-kor-1}

\Aue{Nikiforov, D.} 2015. Ob informatsionnoy bezopasnosti budushchey sistemy situatsionnykh 

tsentrov v~Rossii [About information security of future system of the situational centers in 

Russia]. \textit{Connect} 7--8. Available at: {\sf  

http://www.aladdin-rd.ru/\linebreak company/pressroom/articles/43615} (accessed February~9, 2016). 



\bibitem{16-kor-1}

Strategiya natsional'noy bezopasnosti Rossijskoy Federatsii, utverzhdena Ukazom Prezidenta RF 

ot 31.12.2015 No.\,683, st.\,112 [Strategy of national security of the Russian 

Federation, approved by the Decree of the Russian President of 31.12.2015 No.\,683, Art.\,112].



\end{thebibliography}

} }





\vspace*{-1pt}



\hfill{\small\textit{Received February 9, 2016}}



     

     \Contr

     

     \noindent

     \textbf{Zatsarinny Alexander A.} (b.\ 1951)~--- Doctor of Science in 

technology, professor, Deputy Director, Federal Research Center ``Computer 

Science and Control'' of the Russian Academy of Sciences (FRC CSC);

Institute of Informatics Problems of FRC CSC, 44-2 Vavilov Str., 

Moscow 119333, Russian Federation; \mbox{AZatsarinny@ipiran.ru} 



\vspace*{3pt}



\noindent

\textbf{Korolev Vadim I.} (b.\ 1943)~--- Doctor of Science in technology, leading 

scientist, Institute of Informatics Problems, Federal Research Center ``Computer 

Science and Control'' of the Russian Academy of Sciences, 44-2 Vavilov Str., 

Moscow 119333, Russian Federation; professor, Department ``Cybernetics and 

Information Security,'' National Research Nuclear University ``MEPhI'', 

31~Kashirskoye Sh., Moscow 115409, Russian Federation; 

VKorolev@ipiran.ru   

     

 \label{end\stat}





\renewcommand{\bibname}{\protect\rm Литература}

